% !TeX program = pdflatex
% IMT School for Advanced Studies Lucca beamer template.
% Build with `make` (see the Makefile) or with pdfLaTeX in TeXstudio.
\documentclass[aspectratio=169,11pt]{beamer}

% Theme options: sectionpages=true|false, footer=true|false, assetpath=<dir>.
\usetheme{imtlucca}

% -----------------------------------------------------------------------------
% Content packages
% -----------------------------------------------------------------------------
\usepackage{amsmath,amssymb}
\usepackage{booktabs}
\usepackage{appendixnumberbeamer}
\usepackage{microtype}

% Figures path.
\graphicspath{{assets/}{figures/}}

% -----------------------------------------------------------------------------
% Optional co-institution logo, shown on the title slide next to the IMT logo.
% Point it at your own file, or delete the line for an IMT-only title slide.
% -----------------------------------------------------------------------------
% The two marks balance when both are the horizontal lockup at a matched height.
\imtlogo[11mm]{imt_lucca_horizontal.pdf}
\imtcologo[11mm]{logo-uniud.jpg}

% -----------------------------------------------------------------------------
% Metadata: replace these placeholders.
% The optional arguments are used in the footer: \author[Speaker]{...},
% \title[Short title]{Full title}.
% -----------------------------------------------------------------------------
\title[Short title]{Presentation Title}
\subtitle{Optional subtitle}
\author[Speaker]{\underline{Speaker Name} \and Second Author \and Third Author}
\institute[IMT Lucca]{IMT School for Advanced Studies Lucca}
\date{Venue \\ \today}

% -----------------------------------------------------------------------------
% Document starts here.
% The following frames are generic placeholders only.
% Replace them with your own content.
% -----------------------------------------------------------------------------
\begin{document}

\begin{frame}[plain,noframenumbering]
  \titlepage%
\end{frame}

\begin{frame}{Outline}
  \tableofcontents
\end{frame}

\section{Template test}

\begin{frame}{Clean text slide}
  \begin{itemize}
    \item Replace this placeholder with your own first point.
    \item Use \highlight{orange} for emphasis and \bluehighlight{blue} for secondary emphasis.
    \item Keep slides sparse: one idea per slide works best for a 30-minute talk.
  \end{itemize}
\end{frame}

\begin{frame}{Two-column layout}
  \begin{columns}[T,onlytextwidth]
    \begin{column}{0.48\textwidth}
      \begin{block}{Left block}
        Placeholder text for definitions, assumptions, or setup.
      \end{block}
    \end{column}
    \begin{column}{0.48\textwidth}
      \begin{exampleblock}{Right block}
        Placeholder text for examples, results, or implications.
      \end{exampleblock}
    \end{column}
  \end{columns}
\end{frame}

\begin{frame}{Equation and key message}
  \begin{equation*}
    \text{score}(x) = \alpha \cdot \text{signal}(x) - \beta \cdot \text{cost}(x)
  \end{equation*}

  \vspace{3mm}
  \begin{takeawaybox}
    Replace this sentence with the main conclusion of the slide.
  \end{takeawaybox}
\end{frame}

\begin{frame}{Results formalized in Lean}
  Use \texttt{leanbox} for statements that are proved and implemented in the
  Lean proof assistant. The first argument is the kind of result, the second is
  its name; leave the name empty to drop the parentheses.

  \begin{leanbox}{Theorem}{Soundness}
    If $\Gamma \vdash e : \tau$ then $e$ evaluates to a value of type $\tau$.
  \end{leanbox}

  \begin{leanbox}{Lemma}{}
    Every well-typed term is either a value or takes a step.
  \end{leanbox}
\end{frame}

\begin{frame}{Figure placeholder}
  \begin{center}
    \fbox{\parbox[c][42mm][c]{0.72\textwidth}{\centering Insert figure here}}
  \end{center}
\end{frame}

% Keep each panel to one sentence: four panels leave little vertical room, and
% bullet lists here shrink into something nobody reads from the back row.
\begin{frame}{In summary}
  \vfill
  \begin{recap}
    \recapitem{Problem}{State the gap your work addresses, in one sentence.}
    \recapitem{Methodology}{Name the approach, not its details.}
    \recapitem{Results}{Give the single number or property that matters.}
    \recapitem{Future work}{Say what comes next.}
  \end{recap}
  \vfill
\end{frame}

\begin{frame}{Thank you}
  \Large Questions?

  \vspace{6mm}
  \normalsize
  \textcolor{IMTMutedText}{Name Surname \quad | \quad email@example.org}
\end{frame}

\appendix
\section*{Backup slides}

\begin{frame}{Backup slide placeholder}
  This area is for optional details, proofs, extra experiments, or robustness checks.
\end{frame}

\end{document}
